\documentclass[11pt]{article}

\usepackage[letterpaper,margin=2cm]{geometry}
\usepackage{amsmath,amssymb}

\begin{document}
\thispagestyle{empty}

\begin{center}
    {\Large\textbf{Ejercicio e) - Cambio de variable}}
\end{center}

\[
L=\lim_{x\to 2}\frac{x-\sqrt{x+2}}{\sqrt{4x+1}-3}
\]

\noindent Tomamos el cambio correcto:
\[
\beta=\sqrt{x+2}
\quad\Rightarrow\quad
x=\beta^2-2,
\qquad
x\to 2 \Rightarrow \beta\to 2.
\]

\noindent No se debe poner tambien \(\beta=\sqrt{4x+1}\); esa raiz se
reescribe usando \(x=\beta^2-2\):
\[
\sqrt{4x+1}
=\sqrt{4(\beta^2-2)+1}
=\sqrt{4\beta^2-7}.
\]

Entonces
\[
L=
\lim_{\beta\to 2}
\frac{\beta^2-2-\beta}{\sqrt{4\beta^2-7}-3}
=
\lim_{\beta\to 2}
\frac{(\beta-2)(\beta+1)}{\sqrt{4\beta^2-7}-3}.
\]

Racionalizamos el denominador:
\[
\sqrt{4\beta^2-7}-3
=
\frac{(4\beta^2-7)-9}{\sqrt{4\beta^2-7}+3}
=
\frac{4(\beta-2)(\beta+2)}{\sqrt{4\beta^2-7}+3}.
\]

Por lo tanto,
\[
L=
\lim_{\beta\to 2}
\frac{(\beta-2)(\beta+1)(\sqrt{4\beta^2-7}+3)}
{4(\beta-2)(\beta+2)}
=
\frac{3(3+3)}{4(4)}
=
\boxed{\frac{9}{8}}.
\]

\vspace{0.4cm}
\noindent\textbf{Verificacion con Python:} SymPy entrega \(\frac{9}{8}\).

\end{document}
